\section{蜀汉的战争与行政：从成都到剑阁的每一次转手}

蜀汉的行政不应只在丞相府中寻找。成都发出一封任命、一纸催粮文书，到了汉中，
才会变成一座仓口是否开启、守军是否按时换防；到了南中，才会遇上郡县官与地方
首领能否相互借力的问题；沿峡江向东，则又是船只、戍卒和对吴使者能否通行的
安排。盆地的农业和聚落足以支撑一座朝廷，却没有留下连续的产量账本。因而不能
把“天府”轻率换算成某一年可以北运多少粮食、征发多少人。较能落实的问题是：
在宫廷、汉中守备、南中联络和吴蜀外交之间，有限的人力与物资怎样被一次次改派；
每次派遣又如何改变了下一次选择的余地。

这里的年序主要依《三国志》蜀书及裴松之注建立。《三国志》让人看见诏令、任免、
出军和退军，却很少把一名县吏怎样催集夫役、一户人家怎样承受征调逐日写下。
裴注保存的《汉晋春秋》《襄阳记》等材料，常能把一处营地或一段争论照亮得更近，
但它们成书、立场和叙事目的并不相同，不能自动补成现场记录。《华阳国志》保存
巴蜀与南中的地方传统，尤其值得用来抵抗“成都以外皆为空白”的写法；它同样是
后出的地方史，不能单独为每一座山、每一位首领和每一次行军作即时见证。读者须
在这些材料的交错处，辨认一个小国怎样把战争变为治理，也辨认它不能治理到哪里。

\subsection{永安的遗命与成都的接手}

\begin{event}{托孤之后，谁来把命令送回成都}{章武三年四月（223年）}{蜀汉}{\eventref{TK-0223-LIU-BEI-DEATH}{刘备去世与辅政}}
刘备死于永安，刘禅继位，诸葛亮受遗诏辅政。这个顺序见于《三国志·先主传》
与《三国志·诸葛亮传》；它可以确认最高权力的交接，却不能把白帝城中的谈话、
各方的心理或地方官的即时态度写成可见的事实。
\end{event}

永安在峡江上游，离成都并不只是地图上的一段距离。君主死去以后，诏令、印信、
随从、守军和消息都要沿既有道路回到盆地，才可能把“受遗诏辅政”变成朝廷日常。
《先主传》把刘备临终托付诸葛亮、李严的名义关系保存下来；《李严传》又显示，
这位被委以重任的官员并不只是一行遗命中的陪衬。对刚继位的刘禅而言，最急迫的
并非抽象地选择“信任谁”，而是让成都的文书系统、军队指挥和地方任官不因永安
的消息而各自停摆。

诸葛亮后来集中辅政权力，往往使读者倒过来把二二三年看成一条必然通向丞相独揽
的道路。材料并不支持这种倒读。它能支持的是：刘备死亡使蜀汉必须同时处理东部
防线、成都中枢和年轻君主的名义位置；诸葛亮与李严都进入这一安排。它不能让人
看见每一名官员在遗命到达时已经站定的派别。正因为继承的实际执行有不确定性，
朝廷才格外需要用任命、尚书文书和军令来把不确定性压进可操作的程序。

这一程序也有社会的背面。王朝的继承常在传记里只占几句话，沿途的驿传、护送和
驻守却不由传记逐项登记。永安至成都之间的军民并不是只等候一则宫廷新闻；他们
必须判断谁的命令可以征粮、谁能调兵、原有的事务是否仍然有效。我们没有足够的
地方文书替这些人写出统一反应，但正因此不能把政治交接写成宫门内完成的仪式。
蜀汉此后能够北向，先以这种不显眼的连续性为条件。

\begin{causalchain}[继承不是北伐的序幕，而是行政能力的先决条件]
结构性条件是蜀汉的核心资源集中在成都盆地，而东部与北部又各有军事压力。刘备
在永安去世是触发因素；遗命和辅政任命提供了名义上的接手办法。直接后果是中枢
得以继续发出命令，中期影响则是诸葛亮能够把南中、汉中和对吴关系纳入同一套优先
顺序。这个安排并没有消除地方摩擦；它只是使朝廷有了重新处理摩擦的入口。
\end{causalchain}

\subsection{南中不是一座已经装满的后方仓库}

\begin{event}{南征后的郡县与中介关系}{建兴三年（225年）}{蜀汉与南中地方集团}{\eventref{TK-0225-SHU-NANZHONG}{诸葛亮南征}}
诸葛亮南征，益州郡、永昌一带的叛乱被平定，蜀汉重新安排地方统治关系。
《三国志·诸葛亮传》提供行动与结果的年序；《华阳国志》补入地方脉络。二者均
不足以证明南中各地从此以同一种方式服从，也不足以证实流传甚广的每一段劝降细节。
\end{event}

二二五年以前，成都面对的南方不是一块可以从地图上圈出的同质区域。益州郡、越巂、
牂牁、永昌及其山谷、道路和聚落之间，本已有不同的权力关系。雍闿、高定等人的
活动、孟获在后出叙事中的显著位置，使这段历史很容易被压成“七擒七纵”的单线
故事。《三国志·诸葛亮传》写诸葛亮南征、秋季平定；裴注所引材料和《华阳国志》
留下更多人物与地方传说。它们的价值在于提醒人们：蜀汉的军队进入的不是无人地带。
它们的限度也同样清楚：叙事越具体，越须分辨其距离事件的远近，不能把传奇中的
反复擒纵当作逐次核验过的军情日志。

更可靠的核心在于，军事胜利之后，蜀汉没有简单把所有当地首领逐出。《诸葛亮传》
记其“皆即其渠帅而用之”，并以当地人士处理当地事务。这里可以看见一种受条件
限制的选择：成都若把每一处山地都改由外来官员直接管理，既需更多驻兵，也更难
取得道路、语言、婚姻网络和物产的信息；保留或任用地方中介，则能降低即时占领
成本，却意味着朝廷对地方的控制经由他人实现。所谓“平定”由此不是一把锁落下，
而是一套需要反复维持的关系。

《华阳国志》关于南中的记述尤其适合放在这一场景旁边读。它使永昌、益州郡及地方
人物不只作为成都军报中的地名出现，也保存了地方史家选择记忆的方式。但它成书
较晚，且有把复杂关系归入可传颂人物的倾向。它能帮助提出“成都要凭谁在南中取得
消息和劳力”这个问题，不能替我们统计当地人口、矿产或贡赋，更不能证明每一支
山地群体都把蜀汉视为同一种政治对象。南中人的完整自述没有连续传下来；材料的
缺口必须留在叙事中，而不是用一个整齐的“少数族群”名词盖过去。

南中与北线的关系也不宜写成一条资源管道。《三国志》中的蜀汉朝廷当然希望获得
可用的人力、牲畜和财货，地方安定亦能减少腹地的军事牵制；然而现存材料没有一
套逐年收支簿，能够把南中产出精确折算为祁山前线的粮秣。可以说，南征使成都获得
重新安排南方关系、把一部分地方力量纳入其军事与行政网络的可能；不能说某种物产
从此稳定而无摩擦地供养北伐。道路被雨季、山谷和地方合作切断时，“后方”首先是
一项政治难题，其次才是一张供应表。

\begin{sourcenote}[孟获故事能说明什么，不能说明什么]
裴松之注与《华阳国志》保存了比《三国志》正文更丰满的南中人物和劝服叙事，
后世又由此发展出著名故事。这些材料能说明蜀地记忆如何把武力与怀柔放在同一
问题中；它们不能单独证实战斗次数、行军路线、所有参与者的意见，或把“服心”
解释为可量化、永久的地方服从。写南中时，叙事的丰满不能超过材料的承重。
\end{sourcenote}

\subsection{汉中：一封表文走到山口以后}

\begin{event}{驻汉中与《出师表》}{建兴五年（227年）}{蜀汉}{\eventref{TK-0227-ZHUGE-LIANG-HANZHONG}{驻汉中与《出师表》}}
诸葛亮驻汉中，上《出师表》，为北伐重新组织军政。《三国志·诸葛亮传》保存
相关年序与表文；表文是向后主陈述用人、秩序和北向选择的政治文字，不是全蜀
财政账本，也不是诸葛亮内心的逐字记录。
\end{event}

汉中盆地不是一个天然向北发射军队的台阶。它要先被守住：守军需要轮换，仓口要
有人看管，褒斜、傥骆和通往陇右的道路要有向导、修治与传令。成都把物资送出
盆地，往往不是一次“运输”，而是一连串由不同人接手的动作；任何一处延误，都
会把前线的选择改成等待。正是在这个意义上，《出师表》所说的“益州疲弊”不宜
读成夸张修辞后便搁置，也不宜据此虚构精确数字。它把北向行动置于一个资源有限、
必须谨慎任人的公开理由之下。

表文首先讨论的其实是人：宫中、府中、亲贤、远佞，以及赏罚应否同一标准。这不是
与军粮无关的道德开场。一个准备长期守汉中、再越山出军的政权，需要让县吏、将领、
尚书和近侍相信命令从何处发出、谁有资格更改、失败后由谁受责。诸葛亮把“先帝
创业未半而中道崩殂”置于表文开端，既是在向年轻君主说明局势，也是把二二三年
以来的继承压力重新写入北伐计划。军事行动若没有可被朝廷承认的任用与赏罚，就
很难把远方的临时协作维持成一条供应链。

这封表文也揭示了蜀汉政治的尺度。它给读者看见丞相怎样向皇帝陈述，不给读者看见
每个家庭是否赞成出军。因征调离开田地的人、在仓口服役的人、等候丈夫或儿子归来
的人，几乎没有自己的连续声音。不能因为材料沉默便假设他们全都热情或全都反对；
能做的是把他们放回选择的成本中：每一次在汉中集结，都意味着成都及其周边须把
劳力、牲畜、车具和地方秩序的一部分让给前线。

\subsection{祁山的响应与街亭的一口水}

\begin{event}{祁山出兵与街亭失守}{建兴六年春（228年）}{蜀汉与曹魏}{\eventref{TK-0228-JIETING}{街亭失守}}
诸葛亮出祁山，南安、天水、安定一度响应；街亭失守后，蜀军撤回汉中，并处分
马谡等人。《三国志·诸葛亮传》《马谡传》与魏方纪传可相互校验行动骨架。裴注
中关于布阵、断水的说法增添战场细节，却不应被误作双方都留下的现场笔录。
\end{event}

祁山方向的最初变化，使战争暂时不再只由成都和洛阳的命令决定。南安、天水、安定
的响应，对蜀军重要的不仅是“得三郡”的名声：郡县若转向，军队或许得到可停驻的
城、可询问的地方消息、可征集或购买的粮食，以及能够绕开敌军的道路知识。可这些
条件都带着期限。刚刚响应的郡县未必有足够的兵力自守，也未必愿意把所有存粮交给
远来的军队；蜀军若不能保持交通，原本给它增加纵深的地方关系，反会成为难以承受
的保护责任。

街亭的位置因而不能只作为马谡传中的性格考题来读。《马谡传》记马谡违亮节度、
为张郃所破；《三国志》魏方材料也保留张郃击破马谡的结果。裴注所引《襄阳记》
等说法把山上扎营、取水受阻的画面带到读者眼前。这样的画面能够解释为什么一处
小小的据点会关乎大军退路，却应保留其材料层次：我们可以确认失守和撤军的先后，
不能从一种后出叙事重建每一道壕沟、每一声命令，更不能把一场战败简化为某一个人
“骄傲”或“愚蠢”的自然报应。

诸葛亮撤军并处置马谡、张休、李盛等人，是朝廷在失败后做出的可见动作。处分有
实际的政治功能：它告诉各军，谁对违令和失守负名义责任；它也告诉仍在观望的郡县，
蜀汉并未放弃以军法约束自己的指挥链。但处分不能把道路、时令、魏军援军和地方
合作的复杂关系全部装进一人的罪名。将马谡写成全部原因，会让读者错过更难而更
重要的事实：蜀军没有一条可以轻易替换的安全补给线，一处节点不能维持，前方三郡
的短暂响应便失去可持续的依托。

撤回汉中之后，失败仍留在社会中。退军要带回伤者，清点未能带走的器具，重新部署
守山口的人；被征发而未随军的人则面对下一次征调的可能。史书让我们看见将领被
贬黜，难以让我们逐户看见这些成本怎样分配。正因为不可见，叙事更不应以“蜀人
同仇敌忾”这种总括来替代。北伐不只是军队向北移动，也是在盆地内外重排谁必须
等待、运输、服役和承担风险。

\begin{turningpoint}[街亭之后，朝廷能惩罚失守，却不能惩罚距离]
街亭失守改变的不是一段传记的评价，而是祁山方向的政治地理。三郡的响应需要
蜀军能够停留才有意义；停留又取决于道路、据点、仓储和地方选择。军法可以重建
指挥的可追责性，不能凭空建成一条新的山地补给线。此后取武都、阴平和再出祁山，
都是在这个限度内重新寻找较可承受的位置。
\end{turningpoint}

\subsection{粮道成为责任：李严案与再出祁山}

\begin{event}{祁山撤军与李严被废}{建兴九年（231年）}{蜀汉}{\eventref{TK-0231-LI-YAN-DISMISSAL}{李严被废}}
诸葛亮再出祁山，与魏军相持，因粮运不能接续而撤军；李严随后因转运失职及伪称
诏命被弹劾、废为平民。《三国志·诸葛亮传》《李严传》保存诏令、弹劾和处分的
骨架，不能据此把全部补给困难归结为一名官员。
\end{event}

二三一年再出祁山时，蜀汉已从街亭的失败中知道，北方的机会不能脱离运输。李严
负责督运，正说明后方事务不是“军队出发以前”的准备工作，而是决定军队何时必须
停止的前线条件。《李严传》所记情节有清晰的制度性轮廓：粮运不继，诸葛亮退军；
李严又以诏书名义召诸葛亮回军，事后被揭发而遭弹劾。这里最值得注意的不是把人物
立刻判成忠奸，而是朝廷为何必须将运输故障写成文书责任。

远征中的补给困难通常有许多层次：道路是否可通，雨水与季节如何影响转运，牲畜、
民夫和仓储是否足够，前线是否改变驻地，地方是否能够继续供给。现存蜀汉材料没有
把这些层次逐日列成账目。李严案提供的是一场公开追责的档案，不是一份完整的物流
报告。诸葛亮及朝廷通过弹劾重申：前线的撤退不能任意伪托诏命，承担督运者必须对
文书和结果负责。它使中枢重新获得判断命令真伪的权力，却也把一个系统性的脆弱处
压缩成一名高官的失职。

李严被废的后果不只在名位。他原来所掌握的军政关系、随从和地方联系，都需由其他
人接收；诸葛亮则必须在继续北进、整顿道路和安抚朝廷之间重新计算。对普通承担
转运的人而言，朝廷的处分并不会自动缩短山路。它可能使下一道命令更清楚、使某些
官员更谨慎，却不能保证驮畜不死、雨水不涨、仓口必然有粮。因此，李严案应被读作
行政在压力下创造责任边界的努力，而非一个人挡住或放出了整个北伐的寓言。

这种努力有一个政治代价。一个小国必须让官员相信失败会被追究，才可能维持跨越
成都、汉中和前线的协作；追责若只服务于卸责，则会使愿意报告困难的人更少。我们
无法从现有材料量化李严案之后的官场心态，但可以看见其制度困境：朝廷要得到真实
的运输消息，又要在撤军后向君主、将士和地方证明命令仍然有效。二者不总能彼此
兼得。

\subsection{五丈原：让时间替粮道说话}

\begin{event}{五丈原的屯驻与撤军}{建兴十二年（234年）}{蜀汉与曹魏}{\eventref{TK-0234-WUZHANG}{五丈原屯驻}}
诸葛亮出斜谷屯五丈原，与司马懿相持，病逝后蜀军撤回汉中。《三国志·诸葛亮传》
可确认出军、病逝与撤军的顺序；裴松之注所引《汉晋春秋》关于屯田和渭滨居民的
叙述，照亮驻军与耕作的关系，但不能证实营垒规模、兵粮总数或每次军议。
\end{event}

五丈原的场景不是两位统帅隔着战场等待一场决定天下的会战。蜀军从斜谷出来，在
渭水一带停驻，首先要面对的是停留本身能否延续。裴注引《汉晋春秋》说诸葛亮因
常患粮食不继，分兵屯田，耕者杂于渭滨居民之间，而百姓得以安堵、军中不私取。
这段文字之所以重要，不在于它可以替我们画出一张精确的营田平面图，而在于它把
一场相持放回生产、居民与军纪之间：前线若想久驻，不能只从成都拖运，也要尽量
让一部分消耗在当地得到处理。

“百姓安堵”的说法需要谨慎读。它出自后出史料经裴注保存，带有赞扬诸葛亮军纪的
叙事目的。它可以说明作者认为军队与渭滨居民的关系是评价屯驻成败的重要尺度，
不能证明每一个村落都毫无扰动，更不能据此断言蜀军从不征用、不冲突。军队在陌生
地区耕作、驻守和取水，本身已会改变道路、田地和劳力的分配。材料让人看到一种
治理理想如何被表达，却不能让人跳过战争对地方生活的压力。

魏军采取相持，也不是没有成本的静止。关中防线、援军、侦察和其他方向的调度同样
需要资源；只是魏更接近其关中基础，能让时间首先暴露蜀军来路较长的难题。司马懿
与诸葛亮的选择皆受资源约束。把相持写成一方“畏战”、一方“求战”，会把粮道、
屯田和地理从解释中删掉。能够确认的是诸葛亮病逝、蜀军撤回；关于死后军中如何
传令、营垒怎样掩护、魏军是否立即识破等细节，各种后出叙事彼此不一，不能用最
戏剧化的一种填满材料留下的空白。

诸葛亮死后由杨仪等率军撤退，蒋琬录尚书事。撤军不是一张句号：汉中需要接回
部队，成都需要重新分配官署与军权，留在路上的器物与人员需要被处置。二二三年
以来围绕辅政建立的中枢，到此必须面对一个新的问题：没有诸葛亮在场，蜀汉能否
仍把前线经验转为行政常规？答案不在某一句遗言，而在后来的守备、任命和谨慎的
战略选择中慢慢显露。

\begin{textwindow}[渭滨屯田不是“后勤奇迹”]
裴松之注引《汉晋春秋》称诸葛亮“每患粮不继”，因而分兵屯田，“耕者杂于渭滨
居民之间，而百姓安堵，军无私焉”。这段短文的语境是解释久驻之策，并称颂军纪。
它能照亮蜀军试图缩短供应距离的办法；不能证明屯田规模、粮食产量、所有居民的
感受，亦不能代替《三国志》所提供的出军与撤军年序。
\end{textwindow}

\subsection{蒋琬、费祎：不出大军也是一种安排}

诸葛亮死后的年代常被概括为“守成”，仿佛战争一停，行政便只是在维持旧物。
《三国志·蒋琬传》提供了另一种较具体的画面。蒋琬接掌尚书事、后为大司马，并曾
思考由汉水、沔水方向出兵的可能。朝臣担心水路艰险，尤其担心进退不易；蒋琬身体
也日渐有病。材料没有要求读者把这项计划称作正确或错误，却让人看见：同一条水路
在地图上似乎绕开了部分山道，在实际决策中却带来舟师、回撤、河道与地方接应的新
风险。行政不是把“进攻”换成“保守”，而是在不同风险之间选择较能承受的一种。

\begin{event}{蒋琬接掌中枢}{建兴十二年（234年）}{蜀汉}{\eventref{TK-0234-JIANG-WAN-REGENT}{蒋琬接掌中枢}}
诸葛亮死后，蒋琬录尚书事，继掌中枢政务。《三国志·蒋琬传》可确认这一任命及
其后战略讨论；它不能证明所有官员都以同一理由接受减少大规模北伐的安排。
\end{event}

蒋琬所面对的是退军后的多重收尾。汉中的守军不能因主帅去世而撤尽，成都的官署也
不能把所有资源继续押在一场未完成的远征上；同时，吴蜀联盟仍需使者和边境警戒来
维系。把这些事务写成“没有战功”的空白，会误解小国如何存续。守备、仓储、轮调、
安置与外交不一定留下传奇名字，却使下一次危机来临时还有可调动的人和道路。

费祎的时期将这种选择说得更直白。《三国志·费祎传》记姜维多次请求大举北伐，
费祎认为蜀汉国力不如魏，不宜轻出，许其所领兵不过万人。传记以人物言论保存
这一冲突，未必能让人听见朝廷中每一种意见；但冲突本身很重要。姜维从边地军情
看到机会，费祎从国家承受力看到上限。二人不是一个代表勇气、一个代表胆怯，而是
分别把同一政权的两种时间尺度带进决策：窗口稍纵即逝，人口、粮道和守备却不能
在窗口关闭后立刻复原。

费祎被郭修刺杀后，蜀汉失去的也不只是一名大臣。长期辅政依赖能够在君主、尚书、
将领和地方官之间传递判断的人；一旦这样的人突然死去，既有的平衡便须重新由他人
承担。材料难以量化刺杀如何改变每一项预算或每一个派系，却可以提醒读者，蜀汉的
政治稳定不是无人的制度机器。它靠具体官员维持，而这些官员同样暴露在宫廷、使节、
护卫和私人关系的风险之中。

\begin{debate}[“休养生息”不能抹去汉中和边地]
后人常用“休养生息”概括蒋琬、费祎时期。这个概括若指减少不易承担的大规模
远征，可以帮助理解政策取向；若暗示边防、转运和征调已经停止，便超出材料。
《蒋琬传》《费祎传》让人看见战略争论与任官，不提供一张全蜀逐年赋役表。较稳妥
的说法是：中枢试图把资源保留给守备与可控行动，而这种保留本身也需要地方社会
持续承担。
\end{debate}

\subsection{姜维的机会，及其落在地方上的重量}

\begin{event}{陇西羌胡起事与姜维出兵}{延熙十年（247年）}{蜀汉、曹魏与陇西羌胡集团}{\eventref{TK-0247-JIANG-WEI-LONGXI}{陇西出兵}}
姜维趁陇西羌胡起事出兵，魏将郭淮等应战，蜀汉重新把军事重心推向陇右。
《三国志·姜维传》保存蜀魏接触的骨架。汉文史料主要从军政视角记录羌胡，
不能据此写出各部统一的意图、人口或固定政治边界。
\end{event}

姜维北向并非突然抛开蒋琬、费祎年代的谨慎。正相反，它是在汉中守备、吴蜀联盟
和长期边地接触已经存在的条件下，试图利用陇西局势打开机会。羌胡各部的起事、
迁徙和结盟，对蜀魏双方都可能造成牵制；但他们不是两国任意移动的棋子。现存叙事
往往只在蜀军抵达或魏军应援时才让他们显影，这意味着我们可以确认一些战事与接触，
却不能把每一部落的选择写成蜀汉的附属行动。

对成都而言，机会也有账面之外的重量。远征需要将领、士卒、马匹和粮食从本已紧张
的守备中抽出；在前线得到一段道路或一座城，并不等于能长期驻守。对汉中与盆地的
人而言，频繁的出兵意味着征调的可能持续存在，田地和家庭要为不确定的归期留出
余地。传记没有把这种压力逐户登记，不能替他们假造悲惨或欢欣的统一表情；然而若
只报战绩，便会把战争真正赖以发生的劳力和等待从历史中抹掉。

二六二年狄道之役把这种反作用写得更清楚。姜维围狄道，邓艾据险救援，蜀军在粮运
受压、魏军反击的条件下撤退。这个结果见于《三国志·姜维传》及魏方相关传记，
足以说明机会没有自动变为稳定据点；它不足以支持常见的精确伤亡、粮数和行军日程。
狄道的撤退也不应被孤立地宣判为“灭蜀的唯一原因”。更合适的理解是，一次次未能
转成可持续驻守的北向行动，持续压缩了成都能够同时照顾汉中、南中、东部和中枢的
余裕。

\begin{event}{狄道撤军}{景耀五年（262年）}{蜀汉与曹魏}{\eventref{TK-0262-JIANG-WEI-DIDAO}{狄道撤军}}
姜维围攻狄道，邓艾据险救援，蜀军因粮运压力与魏军反击撤退。核心叙事见
《三国志·姜维传》；兵力、伤亡、每日行程及地方群体的态度没有连续账簿支撑，
不应以后出的战争故事补足。
\end{event}

狄道之后，蜀汉朝廷的争论不只是“还要不要北伐”。它还涉及谁掌握军队、谁向后主
报告边情、谁为汉中留下足够守军，以及地方官怎样应付再次征发。黄皓等宫廷人物在
后世叙事中常被塑成解释一切的符号，但单一近侍不能取代道路、军队和政治程序的
分析。《三国志·后主传》《姜维传》能够让人看见宫廷与将领的矛盾线索，却不足以
把蜀汉末年的全部失败归因于一人的谗言或一名将军的固执。将复杂压力道德化，是
最省力也最不可靠的写法。

\subsection{阴平以后：政权终结不等于地方瞬间换人}

\begin{event}{邓艾入成都与刘禅出降}{景耀六年（263年）}{蜀汉与曹魏}{\eventref{TK-0263-CHENGDU-SURRENDER}{刘禅出降}}
魏军多路攻蜀，钟会在剑阁受阻，邓艾越阴平逼近成都平原，刘禅出降，蜀汉政权
终结。《三国志·后主传》《邓艾传》《钟会传》可互证大致顺序；阴平道路的艰险
与成都的震动不等于每一处细节、每位官员的动机都可确证。
\end{event}

二六三年把长期分散的压力集中到几条道路上。钟会在剑阁方向受阻，说明蜀汉并非
没有依山设防的能力；邓艾改由阴平逼近成都，则表明一套防御若把资源集中于熟知的
关口，仍可能被另一条高成本、难以预料的路线绕过。《邓艾传》以强烈的功业笔法
叙述这次行动，《后主传》则从成都朝廷的角度记录出降。二者共同支持政权迅速终结
的顺序，却不允许把所有山道上的困难、军队数量和成都人的反应写成毫无争议的全景。

刘禅的出降常被后人处理成一句人格判词；从行政史看，它首先是一个极端条件下的
资源判断。成都面对的是逼近的魏军、可否调回的汉中部队、城内官员和居民的安全，
以及继续抵抗是否还能改变结果。我们不知道每一次议论的完整记录，不能替刘禅或
诸臣写出内心独白。可以确认的是，最高政权的名义在出降中终止；不能说蜀地的官吏、
军人、地方中介、南中关系和山道交通从此立即变成一种整齐的魏晋秩序。

接收一个区域世界，需要的远不止一封降表。新政权要接触旧有的官吏名册、仓储、
道路、驻军和地方关系；原先受蜀汉任命的人则要判断新命令能否执行、旧部众如何
安置。姜维在剑阁一线的部队、成都的官署、南中和汉中的不同位置，面对的速度并不
相同。现存史书较偏重胜者的接收与将领功业，难以让人看见普通人如何调整生计；
这种不均衡本身应当成为读者的警觉，而非用“蜀地已定”四字抹平。

蜀汉的战争与行政因此留下一个并不浪漫的结论。它能在成都集中命令，也能通过
汉中、南中中介和山道守备把命令延伸出去；但每一次延伸都会增加运输、任官、地方
合作与社会负担的成本。南征并未把后方变成无穷仓库，北伐也从未只是将领意志的
直线。街亭、李严案、五丈原、蒋琬与费祎的取舍、姜维的再进和阴平的绕行，都是
同一个问题在不同年份的答案：一个资源有限的国家，怎样在不确定的道路上维持自己；
又会在何处因无法同时承担所有方向而失去选择。
